% Chapter 8 -- Discussion and Limitations
% Standalone, compilable draft (FPT-template article class + booktabs), sibling to
% sec_4_results.tex / sec_5_recall_bottleneck.tex / sec_7_results.tex. Synthesises the
% detection and completion results into a single statement of what the thesis supports,
% then states every audit-identified limitation BY NAME so the committee cannot surface
% one the thesis has not already named.
%
% MERGE/NUMBERING NOTE: this file opens with \section{Discussion and Limitations} for
% standalone compilation and numbering, but at final assembly the assembler promotes it
% to a chapter (Chapter 8). Forward/back references to sibling chapters and their tables
% use typed numerals ("Chapter~4", "Section~7.4", "Table~4.1") because the chapter files
% do not share a label namespace when built separately.
%
% Honesty-phrasing rules applied (personal_agenda P1):
%   - claims matched to evidence tier (measured / inferred / speculative) explicitly
%   - "mean IoU" always qualified as "of matched detections"
%   - recall denominator restated (>=10 surviving points, Sec 4.1) wherever recall is cited
%   - completion claims scoped per the T9c outcome (PARTIALLY HOLDS) and to static cars
%   - negative results (#40/#45 etc.) presented as findings, not apologised for; #37 as a
%     metric-design lesson
%   - every number traces to a finding / artifact (cited Finding~\#NN)
%   - the six mandatory named limitations are all present: val-as-test, statics-only
%     accuracy, compact-band d2, empty long band, 23% length-estimate fallback (#41),
%     SemanticKITTI label-quality assumption -- plus the clustering-objectness ceiling
%     (Ch 5) and the offline-runtime scope
%
\documentclass[a4paper,12pt]{article}
\usepackage[english]{babel}
\usepackage[utf8]{inputenc}
\usepackage[T1]{fontenc}
\usepackage{lmodern}
\usepackage{amsmath,amssymb}
\usepackage{booktabs}
\usepackage{float}
\usepackage{siunitx}
\usepackage{microtype}
\usepackage[margin=2.5cm]{geometry}
\usepackage{graphicx}
\usepackage{multirow}

\begin{document}

\section{Discussion and Limitations}
\label{ch:discussion}

The preceding chapters reported results; this one interprets them. It states, in one
place, what the evidence in this thesis does and does not support, and then names every
limitation the work carries --- including the ones an examiner is most likely to raise.
The organising principle is deliberate: a limitation that the thesis names, quantifies,
and bounds is a defended limitation; one that a committee surfaces first is not. Every
weakness discussed here is already measured elsewhere in the document, and is
cross-referenced to the finding that quantifies it. Nothing in this chapter softens a
result reported earlier; it draws the boundary around each result precisely.

We first summarise what the thesis establishes and how far each claim reaches
(Section~\ref{sec:disc-what-holds}). We then work through the limitations in four
groups: threats to the validity of the evaluation itself
(Section~\ref{sec:disc-eval-validity}), the scope of the completion claims
(Section~\ref{sec:disc-completion-scope}), the ground-truth and dataset assumptions
(Section~\ref{sec:disc-gt}), and the architectural and operational limits of the
approach (Section~\ref{sec:disc-architectural}). Section~\ref{sec:disc-generalises}
closes by separating what generalises from what is specific to the sequences studied,
which sets the agenda for the future work of Chapter~9.

\subsection{What the thesis establishes}
\label{sec:disc-what-holds}

Two contributions carry the strongest evidence. The first is methodological: the
donor-frame occluded-side coverage metric (Chapter~7). We are not aware of a prior metric
that scores a real LiDAR car's completion against donor frames of the same static
instance --- a reference-leakage-free construction, and it comes with the demonstration
that the obvious alternative --- Chamfer distance
against accumulated pseudo-ground-truth --- is invalid, because accumulated LiDAR is
itself one-sided and so rewards under-completion (the raw partial scored the lowest
Chamfer distance on every real example; Finding~\#26). This is a measured negative
result about metrics, not an opinion. The second is empirical and rests on that metric:
completion adds genuinely unseen surface (median coverage of never-observed surface
within \SI{0.1}{\metre} of \num{0.304} on sequence~08, seven times the symmetry-mirror
baseline of \num{0.043} and against a raw-partial baseline of \num{0.000};
Finding~\#32) and improves the amodal bounding box --- the shipped completed box reaches a median
bird's-eye-view IoU of \num{0.771} against the amodal ground-truth box, up from
\num{0.725} for the raw partial, with the gains concentrated on width, height, and
centre where the sensor saw least (Finding~\#36) --- and both headline claims survive a
pre-registered held-out replication, with named coverage caveats
(Finding~\#42, outcome \textsc{partially holds}; Section~\ref{sec:disc-eval-validity}).

The detection pipeline is a strong engineering result rather than a claim of
state-of-the-art accuracy. It reaches precision \num{0.905}, recall \num{0.730}, and
$F_1$~\num{0.808} on the full sequence~08 (mean IoU of matched detections \num{0.912};
recall denominator: ground-truth car instances with at least \num{10} points surviving
preprocessing, Section~4.1), and it does so with an interpretable, near-annotation-free
clustering front end on modest hardware. Its recall shortfall is not mysterious: it is
root-caused to HDBSCAN splitting large, close cars into fragments (Chapter~5), and that
diagnosis is itself a contribution. Two secondary findings round out the account: the
recall bottleneck is clustering fragmentation with a recoverable-ceiling quantification
(Findings~\#23/\#34), and the sim-to-real classifier gap is total and symmetric, which
makes synthetic pretraining redundant once \num{420000} real clusters are available
(Findings~\#25/\#30) --- a negative result that replaces, rather than supports, the
two-stage-training contribution proposed at the outset.

The rest of this chapter is about how far each of these reaches.

\subsection{Threats to the validity of the evaluation}
\label{sec:disc-eval-validity}

\paragraph{Validation-as-test for detection.}
The most important structural caveat is that the detection numbers are not reported on a
held-out test set in the machine-learning sense. SemanticKITTI provides labels only for
sequences \numrange{00}{10}; the official test sequences \numrange{11}{21} are
unlabelled, so no labelled sequence can serve as a true detection test set. Sequence~08
is the classifier's model-selection (validation) set, and it is also the headline
detection sequence. Its detection recall therefore reads optimistically: the classifier
was chosen partly by its behaviour on this sequence. We disclose this rather than obscure
it, and we bound its impact rather than dismiss it. The exposure is real on both precision
and recall, not precision alone: the classifier sets the precision and also costs recall,
since a genuine car cluster it labels not-car becomes a false negative, and that
acceptance rate was tuned partly on this sequence. The stage decomposition of Section~4.3
bounds that recall cost --- geometric clustering alone reaches recall \num{0.775}, the
classifier lowers it to \num{0.677}, and the track filter recovers it to \num{0.730}
(Findings~\#47/\#49) --- so the classifier's realized recall contribution, which model
selection could optimistically inflate, is confined to that band and is not the bulk of
the shortfall. The dominant recall limit is the upstream clustering ceiling (\num{0.775}
before the classifier acts): it is imposed by clustering, is independent of the classifier,
and holds on sequence~00 as well (Chapter~5), so it does not depend on the seq-08 selection
at all. The validation-as-test structure therefore inflates the precision and the
classifier's few-point recall contribution, not the clustering-imposed ceiling that
accounts for most of the missed cars; we claim the exposure is bounded, not nil, and that
it does not reach the chapter's main recall finding. A distinct caveat is breadth rather
than optimism: the detection headline rests on a single sequence, so per-sequence recall
variation across the labelled set is characterised only in aggregate (Section~4.4), not
sequence-by-sequence. The completion side does not share the validation-as-test weakness: its headline claims are
re-tested on sequence~00 under a pre-registration committed before any sequence-00
evaluation was run (Finding~\#42), and the completion model was trained on synthetic data
and never on any evaluation sequence. The detection and completion contributions thus
carry different evidential weights, and we do not claim a held-out test result for
detection.

\paragraph{Coverage gaps in the held-out completion replication.}
The pre-registered replication returned the outcome \textsc{partially holds}, not
\textsc{holds}, and the reasons are specific and named. Both primary metrics were
decisive wins on the held-out sequence~00 (BEV IoU \num{0.739}~$\rightarrow$~\num{0.766},
$p=\num{1.6e-3}$; donor coverage at \SI{0.1}{\metre} \num{0.000}~$\rightarrow$~\num{0.413},
$p\approx0$; Finding~\#42), so the value of completion does replicate. Two gaps prevent a
full \textsc{holds}. First, the \emph{long band is empty}: no car of at least
\SI{4.6}{\metre} survived the well-observed-static-car construction on sequence~00
(bands: compact \num{17}, normal \num{28}, long \num{0}), so the per-car length estimate's
behaviour on long cars is untestable on the held-out sequence. This is a coverage gap in
the replication, not a weak result --- the sequence~08 long-band evidence stands --- and
the downgrade to \textsc{partially holds} follows the pre-registered taxonomy's
``caveat named'' clause rather than a metric failure. Second, the compact-band
hallucination guard fails on the held-out sequence --- a real failure, not only a naming
convention --- which the next section treats both as a metric-design point and as a
completion limitation in its own right. And the replication rests on a single held-out
sequence: sequence~00 is the one labelled sequence held out from the completion model
(which was trained on synthetic data, never on a labelled sequence), so the completion
value is replicated once, not across a range of sequences.

\subsection{Scope of the completion claims}
\label{sec:disc-completion-scope}

\paragraph{The accuracy claims are for static cars.}
The two quantitative completion metrics --- donor coverage and amodal-box overlap ---
both require a static car by construction. Donor coverage scores completion of one
frame's cluster against other frames of the \emph{same} static car, which only exists
if the car does not move; the amodal box is built from accumulated static-car points.
Moving cars therefore cannot be scored by either metric. We report them separately and
only as a plausibility check: on the shipped sequence-08 output, moving car tracks
complete into plausible car boxes at essentially the same rate as static ones
(\SI{57.9}{\percent}, \num{11}/\num{19}, versus \SI{53.7}{\percent}, \num{225}/\num{419};
Finding~\#44). This supports the statement that motion does not \emph{degrade}
completion plausibility --- expected, since completion is single-frame and carries no
motion smear --- but it is a plausibility parity result, not an accuracy result. We never
claim mover completion \emph{accuracy}; the metrics that would license that claim do not
apply to movers.

\paragraph{The per-car length estimate falls back too often.}
The shipped length geometry uses a per-car length estimate (the \num{90}th percentile of
gate-passed longitudinal extents over a track, plus \SI{0.12}{\metre}), which fixed the
compact-overshoot problem of the earlier fixed prior (compact signed length error
\num{+0.295}~$\rightarrow$~\num{+0.063}~m; Finding~\#36). But it falls back to the fixed
population prior of \SI{4.14}{\metre} on tracks with fewer than five gate-passed frames,
and that fallback fires on \SI{23.0}{\percent} of shipped completed tracks
(\num{119}/\num{518} on the regenerated sequence-08 output; Finding~\#41). This is well
above a comfortable threshold and is not solved: on those tracks the length geometry is
the population prior, not a per-car estimate, and it reintroduces the compact-overshoot
risk the per-car estimate was built to remove. We measured the fallback frequency rather
than assuming it small, and we carry it as a live limitation.

\paragraph{The residual under-extension is a characterised, length-dependent limit.}
Even with an unbiased length estimate, completed clouds remain slightly too short at the
far end, and the amount needed to correct this scales with car length. We tried to remove
it: a pre-registered intervention (decoupling the output radius from the length push, plus
a synthetic-calibrated length-axis fill factor) fixed the target on normal and long cars
but over-extended compacts, failing the pre-registered compact non-regression guard on
both sequences (Finding~\#45). We report this as a finding: the compensation is
irreducibly length-dependent, so a single global fix cannot remove it without regressing
another band, and pursuing per-length compensation would need more cars than the data
safely supports. The far end is therefore the weakest-covered region of the completed car,
and we say so.

\paragraph{The completion model never saw real LiDAR.}
The completion network was trained only on synthetic KITTI-like partial renders. That is
what makes its evaluation leakage-free (Section~\ref{sec:disc-eval-validity}), but it is
also a limitation, and we name it as the completion counterpart of the classifier
sim-to-real gap of Section~\ref{sec:disc-what-holds}. The model never saw real sensor
statistics, and Finding~\#26 attributes the residual flatness of real completions on the
sparsest, most one-sided inputs to exactly this remaining domain gap (alongside imperfect
centre estimation). Closing it would call for real-data fine-tuning, which is itself a
documented negative-result line in this project (Findings~\#16/\#17/\#19) and is left to
future work; the synthetic-only training is thus at once the source of the metric's
validity and a ceiling on real-data crispness.

\paragraph{Compacts are where completion strays, and the guard was band-blind to it.}
This limitation has two faces --- a completion defect and a metric-design lesson --- and
they must be kept distinct. The completion defect: the shipped completion sprays points
outside the true box on \emph{compact} cars. The metric lesson: the original out-of-box
hallucination guard was a pooled median across all cars, and that pooled statistic was
blind to exactly this band-localised failure. The fixed \SI{4.14}{\metre} length prior
that Finding~\#35 sized --- using that same pooled guard --- hallucinates on compact cars
at roughly \num{100} times the compact band's mirrored baseline (\num{0.0433} versus
\num{0.0004}), a failure the pooled number could not see because the other \num{32} of
\num{40} cars sit near it (Finding~\#37). The shipped per-car length estimate, which
replaced that fixed prior, reduces the compact hallucination to about \num{16} times the
baseline (\num{0.0065}) --- roughly a sevenfold reduction, not elimination --- but on the
\SI{23}{\percent} of tracks that fall back to the fixed \SI{4.14}{\metre} prior (the
fallback of the previous paragraph; Finding~\#41) the full \num{100}-times compact
hallucination returns. So the defect is real, is worst on the fallback tracks, and is not
eliminated on any track.

We report ratios rather than a pass/fail bit because the guard bar is strict and no
shipped configuration clears it on compacts. The bar is that the out-of-box fraction must
not exceed the band's mirrored baseline, so any ratio above \num{1} is a failure and
higher is worse. By that bar sequence~08 fails worst (about \num{16} times with the
shipped estimate, \num{100} times on the fallback tracks), and the held-out sequence~00
confirms the defect does not clear off the tuning sequence either --- a milder but still
failing \num{1.8} times (completed \num{0.009} versus mirrored \num{0.005}). The compact
mirrored baseline is itself near \num{0.0004} --- mirroring a compact car barely leaves its
box --- so that band's bar is effectively ``do not leave the box at all,'' and no method
reaches it. We therefore do not claim the guard ``passes'' on compacts; we report the
ratios and let them stand, because tuning the metric until it passes would defeat the
purpose of a hallucination guard. The per-band guard is what makes the defect visible
where a pooled median hid it (Finding~\#37) --- but the underlying fact is a completion
accuracy limitation on compact cars, not merely a property of the metric.

\subsection{Ground-truth and dataset assumptions}
\label{sec:disc-gt}

\paragraph{The amodal ground truth is constructed, and cannot be externally validated.}
The completion evaluation compares against an amodal (full-extent) box built by
accumulating a static car's points across frames, because no external amodal reference
exists for these sequences. We attempted the obvious cross-check --- KITTI raw tracklets
for the drive corresponding to odometry sequence~08 --- and established that it is
impossible: that drive (\texttt{2011\_09\_30\_drive\_0028}) has no tracklet annotations on
the official archive (verified against a control drive that does), and neither do the
neighbouring candidate drives. The amodal ground truth therefore rests on three defences
rather than an external reference: the paired evaluation design tolerates reference noise
because it compares raw versus completed against the \emph{same} box; the construction
guards enforce viewpoint coverage and near-zero motion; and the dimension distribution is
sane (median L~\num{4.14}~/~W~\num{1.75}~/~H~\num{1.47}~\si{\metre}, consistent with
published car statistics). We do not claim the amodal ground truth is error-free; we claim
the paired design is robust to the errors it can plausibly contain.

\paragraph{Recall against all annotated cars, and the two ``$\geq$10-point'' rules.}
The headline recall denominator counts ground-truth car instances with at least \num{10}
points surviving preprocessing (z-filter, voxelisation, denoising, ground removal). This
excludes cars that are too sparse or too distant to survive the front end, so the headline
recall is a recall against \emph{eligible} cars, not against all annotated cars. We
quantify the gap rather than leave it implicit: on a stride-20 sample of sequence~08,
\SI{25.5}{\percent} of all annotated cars (and \SI{14.6}{\percent} of cars with at least
\num{10} \emph{raw} points) are excluded by the eligibility rule, so recall against all
annotated cars is approximately \num{0.54} (Finding~\#48). The two ``$\geq$10-point''
thresholds are different quantities --- one counts raw label points, the other counts
points surviving preprocessing --- and they are not interchangeable; we keep them
distinct wherever both appear. Stating the \num{0.54} figure alongside the \num{0.730}
eligible-car recall is the honest way to report the operating range: the pipeline detects
about three eligible cars in four, and about half of all annotated cars including those it
never receives enough points to see.

\paragraph{Label quality is assumed.}
All ground-truth counts, matches, and the amodal construction assume SemanticKITTI's
semantic and instance labels are correct. We did not audit label quality, and mislabelled
or mis-instanced cars would propagate into both the detection metrics and the amodal
ground truth. This is a standard assumption for work on this dataset, but it is an
assumption, and we name it.

\subsection{Architectural and operational limits}
\label{sec:disc-architectural}

\paragraph{Density clustering has no notion of objectness.}
The recall ceiling is architectural, not a tuning failure (Chapter~5). Coarsening the
clustering voxel to \SI{0.10}{\metre} recovered part of the loss
(recall \num{0.699}~$\rightarrow$~\num{0.730}; Finding~\#34) by preventing intra-car
splits, but coarsening further would begin fusing adjacent cars, and the
clustering-method benchmark shows that no fixed-radius density clusterer reaches even
HDBSCAN's recall at any tested radius (Finding~\#43). The residual limit is that
density-based clustering groups points by proximity and density alone, with no model of
what a car is, so it cannot be tuned to the recall of a learned object detector. Lifting
this ceiling would require replacing the clustering stage with a learned detector, which
is a different system and is left to future work (Chapter~9). The value of Chapter~5 is to
have located the ceiling precisely --- clustering, splitting rather than
classification-or-merging, and lower than first claimed --- rather than to have removed it.

\paragraph{The system is offline by design.}
The pipeline runs at about \SI{623.5}{\milli\second} per frame on the reference hardware
(\num{1.6}~frames/s; promoted-configuration full-sequence timing, Section~4.5), dominated
by the classical CPU stages (HDBSCAN alone is roughly \SI{59}{\percent}). This is offline,
not real-time, and real-time operation is explicitly out of scope. The runtime work that
was done (Findings~\#33/\#34) reached this operating point and established that the
$\leq\SI{400}{\milli\second}$ target is not reachable through the authorised
exact-output and voxel-resolution tiers, because the post-ground object cloud is
intrinsically sparse; GPU clustering was deliberately not pursued. A real-time variant is
future work, not a claim of this thesis.

\paragraph{Tracking is used but not quantified.}
The centroid tracker is evaluated only indirectly. It feeds the per-car length estimate
and the completion aggregation, and its five-gate-passed-frame requirement is what drives
the \SI{23}{\percent} length-estimate fallback (Finding~\#41): identity switches or
fragmented tracks raise that fallback rate directly, and thereby the compact-overshoot
risk it carries. Yet we report no multi-object-tracking accuracy metric (no MOTA, IDF1, or
equivalent) --- the tracker is used as a filter and an aggregator, not assessed as a
tracker --- so its contribution to the fallback rate, and any completion error it
propagates, is not separately measured. Quantifying tracking robustness is future work.

\subsection{What generalises, and what is sequence-specific}
\label{sec:disc-generalises}

It is worth separating the results that should transfer beyond the two sequences studied
from those that are tied to them. The \emph{mechanisms} generalise, though not all with
the same evidential force, and it is worth grading them. Two transfer by argument rather
than by sample: the demonstration that accumulated-pseudo-ground-truth Chamfer is invalid
follows deductively from the one-sidedness of accumulated LiDAR, which is a property of the
sensor and not of any sequence; and the donor metric's construction is dataset-agnostic by
definition wherever static cars are observed from multiple frames. The third --- the recall
diagnosis that density-based clustering splits large, close, dense cars --- is weaker: it
is an empirical, inductive result, measured on the two sequences studied. It is consistent
across both, and its cause (HDBSCAN cutting along internal density variation within a dense
return) is mechanistic, so we expect it to be a property of the clustering algorithm and
the sensor geometry rather than of sequence~08 --- but that is an inference from two
sequences, not a proven invariant, and we state it as one. The \emph{constants} are more
local. The length prior (\SI{4.14}{\metre}), the length-estimate
quantile and offset, and the geometric-filter thresholds were tuned on sequence~08, and
while the held-out replication shows the completion \emph{value} transfers, the specific
constants are sequence-derived and should be re-checked before applying the system to a new
sensor or a new class. The completion accuracy claims are scoped to static, gate-passed
cars in the compact and normal length bands; the long band is supported on sequence~08 but
was untestable on the held-out sequence, and moving cars are covered only by the
plausibility check. Within that scope the claims are pre-registered and replicated; outside
it, they are hypotheses for future work rather than results.

This is the boundary the thesis draws. The contributions in
Section~\ref{sec:disc-what-holds} hold within it, each limitation above marks where it
ends, and Chapter~9 turns the most consequential of those boundaries --- the clustering
objectness ceiling, the length-estimate fallback rate, and the missing long-band and
mover accuracy evidence --- into concrete future work.

\end{document}
